\documentclass[10pt,a4paper]{amsart}
\usepackage[margin=19mm]{geometry}
\usepackage{amsmath,amssymb,mathtools}
\usepackage[scr=rsfs,cal=euler]{mathalfa}
\usepackage{xfrac}
\usepackage[unicode,hidelinks,bookmarks=false]{hyperref}
\newcommand{\ie}{i.e.\ }
\newcommand{\cf}{cf.\ }
\newcommand{\resp}{respectively\ }
\newcommand{\Subsection}{\subsection}
\providecommand{\nobreakdash}{\nobreak\hbox{-}}
\setlength{\emergencystretch}{2em}
\multlinegap0pt
\usepackage{amssymb}
\usepackage{float}
\usepackage{tikz}
\newtheorem{theorem}{Theorem}
\newtheorem{corollary}{Corollary}
\title{\large \textbf{ISOPERIMETRY BY STRETCHING}}
\author{\small KOBE MARSHALL-STEVENS \& GONGPING NIU}
\date{Source: arXiv:2603.12363v2 (CC BY 4.0)}
\begin{document}
\maketitle

\noindent\emph{Source and licence.} \large \textbf{ISOPERIMETRY BY STRETCHING}. Version arXiv:2603.12363v2 (CC BY 4.0), distributed by arXiv under the Creative Commons Attribution 4.0 International licence, http://creativecommons.org/licenses/by/4.0/, as recorded in the arXiv licence metadata for this exact version and shown on the abstract page https://arxiv.org/abs/2603.12363v2. Full licence notice: \texttt{licenses/arxiv-2603.12363-CC-BY-4.0.txt}.\par\smallskip

\noindent\textit{Editorial scope.} Counted unit: the article's corollary vcm (source0.tex lines 433--444). The article's complete original proof of it is delivered with it (source0.tex lines 446--455, one \texttt{proof} environment of 10 lines). No mathematics is written, repaired or re-derived: the statement and the proof are the article's own lines, in the article's own notation, and no retained line is substituted or re-flowed.  Retained uncounted as the source-local context the proof needs: lines 197--199.  Nothing was omitted from any retained span, so the package carries no omission record.  No retained line contains a citation, so the wrapper carries no bibliography and no bibliography entry.  No retained line contains a figure environment, an image-inclusion command or drawing code, so no image is delivered and the package declares no asset dependency.  Editorial formatting only, in addition: an AMS \texttt{amsart} preamble that restates the article's own declarations and macros used by the retained lines, namely the environments \texttt{theorem}, \texttt{corollary} on separate counters, together with the standard packages those lines need.  The retained environments print in this document as Theorem 1, Corollary 1 in order of appearance, whereas the article prints the same items with the numbering of its own full document.  The complete native source of the article as distributed in the arXiv e-print for this exact version is bundled unchanged as \texttt{sources/196271/source0.tex}.
\begin{theorem}\label{thm: stretching}
    If $\Sigma \subset M$ is a smooth, closed, connected, separating hypersurface which is uniquely homologically area-minimising in a tubular neighbourhood of $\Sigma$ in $(M,g)$, then there is a Riemannian metric, $h$, on $M$ such that $\Sigma$ bounds an isoperimetric region in $(M,h)$.
\end{theorem}

\begin{corollary}\label{cor: stretching for vcm}  
    Suppose that $\Sigma,\widetilde{\Sigma} \subset M$ are smooth, closed, connected, separating hypersurfaces satisfying the following assumptions:
    
    \begin{itemize}
        \item $\Sigma$ is homologically area-minimising in a tubular neighbourhood, $V$, of $\Sigma$ in $(M,g)$.

        \item $\widetilde{\Sigma}$ bounds a unique volume-constrained minimiser in $V$ with respect to the metric $g$ (up to taking complements).

        \item $\widetilde{\Sigma}$ is both homologous to $\Sigma$ and sufficiently close in area to $\Sigma$ with respect to the metric $g$. 
    \end{itemize}
    Then, there is a Riemannian metric, $h$, on $M$ such that $\widetilde{\Sigma}$ bounds an isoperimetric region in $(M,h)$.
\end{corollary}

\begin{proof}
    Denoting $\widetilde{\Omega} \subset M$ an open set bounded by $\widetilde{\Sigma}$, one can follow the construction in the proof of Theorem \ref{thm: stretching} verbatim to ensure that, for $R > 1$ sufficiently large, any isoperimetric region, $E \subset M$, of enclosed volume $\mathrm{Vol}_{g_R}(\widetilde{\Omega})$ in $(M,g_R)$ has at least one connected boundary component, say $S \subset \partial E$, homologous to $\Sigma$ in $V$.
    
    \bigskip
    
    Under the assumption that $\widetilde{\Sigma}$ is sufficiently close in area to $\Sigma$, namely provided
    $\mathrm{Per}_{g_R}(\widetilde{\Omega}) < \mathrm{Per}_{g_R}(\Omega) + \delta$ (where $\partial \Omega = \Sigma$ and $\delta > 0$ is the uniform lower area bound as in the proof of Theorem \ref{thm: stretching}), we now deduce that $E$ must have exactly one boundary component (i.e.~that $\partial E = S$). To see this, since $\Sigma$ is locally homologically area-minimising in $V$ we have $|S|_{g_R} \geq \mathrm{Per}_{g_R}(\Omega)$, and as $E$ is isoperimetric we have 
    $$\mathrm{Per}_{g_R}(\Omega) \leq |S|_{g_R} \leq \mathrm{Per}_{g_R}(E) \leq \mathrm{Per}_{g_R}(\widetilde{\Omega}) < \mathrm{Per}_{g_R}(\Omega) + \delta;$$
    hence $E$ has exactly one boundary component as any additional boundary components would contribute at least $\delta$ perimeter, forcing $\mathrm{Per}_{g_R}(E) \geq \mathrm{Per}_{g_R}(\Omega) + \delta$, a contradiction. The volume-constrained minimisation of $\widetilde{\Omega}$ in $V$ then ensures that $\mathrm{Per}_{g_R}(\widetilde{\Omega}) = \mathrm{Per}_{g_R}(E)$, and the uniqueness implies that $\partial E = \widetilde{\Sigma}$ as desired.
\end{proof}

\end{document}
